% =========================================================================
\section{Motivación: la analogía con las ecuaciones en diferencias}\label{sec:motivacion}
% =========================================================================

Las ecuaciones diferenciales lineales con coeficientes \emph{constantes} se
resuelven con el ansatz exponencial $y=\eul^{\lambda x}$, que transforma la
ecuación en un polinomio característico. Existe un mecanismo idéntico para
las ecuaciones \emph{discretas}, y es precisamente ese mundo discreto el que
reaparecerá, de forma inesperada, cuando busquemos soluciones en serie de
ecuaciones con coeficientes \emph{variables}.

% -------------------------------------------------------------------------
\subsection{Ecuaciones en diferencias con coeficientes constantes}\label{sec:diferencias}
% -------------------------------------------------------------------------

\begin{definicion}[Ecuación en diferencias lineal]\label{def:diferencias}
Sean $k\in\N$, $k\ge 1$, y $a_0,\dots,a_k\in\C$ con $a_k\neq 0$ y $a_0\neq0$.
Una \emph{ecuación en diferencias lineal homogénea de orden $k$ con
coeficientes constantes} es una relación
\begin{equation}\label{eq:diferencias}
  a_k\,f(n+k)+a_{k-1}\,f(n+k-1)+\dots+a_1\,f(n+1)+a_0\,f(n)=0,
  \qquad n\in\N ,
\end{equation}
cuya incógnita es una sucesión $f\colon\N\to\C$. Su \emph{polinomio
característico} es $\chi(\lambda)=a_k\lambda^k+a_{k-1}\lambda^{k-1}+\dots+a_0$.
\end{definicion}

Como $a_k\neq 0$, la ecuación \eqref{eq:diferencias} permite despejar
$f(n+k)$ en función de los $k$ valores anteriores: fijados
$f(0),\dots,f(k-1)$, la sucesión queda determinada por inducción. El conjunto
de soluciones es, por tanto, un espacio vectorial de dimensión $k$, en
perfecto paralelismo con el teorema de existencia y unicidad de las EDO
lineales.

\begin{proposicion}[Ansatz exponencial discreto]\label{prop:ansatz-discreto}
La sucesión $f(n)=\lambda^n$, con $\lambda\neq0$, es solución de
\eqref{eq:diferencias} si y solo si $\chi(\lambda)=0$. Si $\chi$ tiene $k$
raíces distintas $\lambda_1,\dots,\lambda_k$, la solución general es
\begin{equation}\label{eq:sol-diferencias}
  f(n)=C_1\lambda_1^{\,n}+C_2\lambda_2^{\,n}+\dots+C_k\lambda_k^{\,n},
  \qquad C_1,\dots,C_k\in\C .
\end{equation}
Si $\lambda_j$ es raíz de multiplicidad $m_j$, las sucesiones
$n^{i}\lambda_j^{\,n}$, $0\le i<m_j$, sustituyen a las potencias repetidas.
\end{proposicion}

\begin{proof}
Sustituyendo $f(n)=\lambda^n$ y extrayendo el factor común $\lambda^n\neq0$,
\begin{equation}\label{eq:caracteristico}
  \lambda^n\bigl(a_k\lambda^k+a_{k-1}\lambda^{k-1}+\dots+a_0\bigr)=\lambda^n\chi(\lambda)=0
  \iff \chi(\lambda)=0 .
\end{equation}
Por linealidad, \eqref{eq:sol-diferencias} es solución. Para ver que es la
general basta comprobar que se pueden imponer datos iniciales arbitrarios
$f(0),\dots,f(k-1)$: el sistema lineal resultante para $(C_1,\dots,C_k)$
tiene por matriz la de Vandermonde $\bigl(\lambda_j^{\,i}\bigr)_{0\le i<k,\,1\le j\le k}$,
cuyo determinante $\prod_{i<j}(\lambda_j-\lambda_i)$ no se anula por ser las
raíces distintas. El caso de raíces múltiples es análogo al de las EDO y
se omite (véase \cite[cap.~2]{ELAYDI}).
\end{proof}

\begin{ejemplo}[Sucesión de Fibonacci]\label{ej:fibonacci}
La ecuación $f(n+2)=f(n+1)+f(n)$, con $f(0)=0$, $f(1)=1$, tiene polinomio
característico $\lambda^2-\lambda-1$, de raíces
$\varphi=\tfrac{1+\sqrt5}{2}$ y $\psi=\tfrac{1-\sqrt5}{2}$. Imponiendo los
datos iniciales en $f(n)=C_1\varphi^n+C_2\psi^n$ se obtiene la fórmula de
Binet,
\[
  f(n)=\frac{\varphi^{n}-\psi^{n}}{\sqrt5},
\]
exactamente como se resolvería $y''-y'-y=0$ con $y=\eul^{\lambda x}$.
\end{ejemplo}

El paralelismo es completo y se resume en la \cref{tab:analogia}: el
operador de derivación $D=\dif/\dif x$ desempeña en el continuo el papel del
operador de desplazamiento $E\colon f(n)\mapsto f(n+1)$ en el discreto.

\begin{table}[ht]
\centering
\caption{Diccionario entre las EDO lineales con coeficientes constantes y
las ecuaciones en diferencias lineales con coeficientes constantes.}
\label{tab:analogia}
\renewcommand{\arraystretch}{1.35}
\begin{tabular}{@{}l>{\centering\arraybackslash}p{4.6cm}>{\centering\arraybackslash}p{4.6cm}@{}}
\toprule
 & \textcolor{cian}{\bfseries Continuo (EDO)} & \textcolor{cian}{\bfseries Discreto (diferencias)}\\
\midrule
Incógnita            & $y(x)$, $x\in\R$                     & $f(n)$, $n\in\N$\\
Operador básico      & $D\colon y\mapsto y'$                & $E\colon f(n)\mapsto f(n+1)$\\
Ecuación             & $\sum_{j} a_j\,D^{j}y=0$             & $\sum_{j} a_j\,E^{j}f=0$\\
Ansatz               & $y=\eul^{\lambda x}$                 & $f(n)=\lambda^{n}$\\
Polinomio característico & $\sum_j a_j\lambda^j=0$          & $\sum_j a_j\lambda^j=0$\\
Raíz de multiplicidad $m$ & $x^{i}\eul^{\lambda x}$, $i<m$  & $n^{i}\lambda^{n}$, $i<m$\\
Datos que fijan la solución & $y(x_0),\dots,y^{(k-1)}(x_0)$ & $f(0),\dots,f(k-1)$\\
\bottomrule
\end{tabular}
\end{table}

\begin{observacion}[$E=\eul^{D}$]\label{obs:E-expD}
El diccionario no es una mera coincidencia formal. Si $y$ es una función
cuya serie de Taylor en $x$ converge a $y$ en un radio mayor que $1$,
\[
  y(x+1)=\sum_{k=0}^{\infty}\frac{y^{(k)}(x)}{k!}
        =\Bigl(\sum_{k=0}^{\infty}\frac{D^{k}}{k!}\Bigr)y(x)
        \eqqcolon \eul^{D}y(x),
\]
de modo que $E=\eul^{D}$ y el ansatz $\lambda^n=\eul^{n\ln\lambda}$ es el
ansatz exponencial evaluado en los enteros. Esta identidad exige que $y$ sea
\emph{analítica} (\cref{def:analitica}): la noción central de este tema
aparece ya en la motivación.
\end{observacion}

% -------------------------------------------------------------------------
\subsection{Por qué nos interesa: de la EDO a la recurrencia}\label{sec:por-que}
% -------------------------------------------------------------------------

Cuando los coeficientes de la EDO dependen de $x$, el ansatz
$y=\eul^{\lambda x}$ deja de funcionar. La estrategia de este tema consiste
en proponer una \emph{serie de potencias} $y=\sum a_n(x-x_0)^n$ y traducir la
ecuación diferencial para $y(x)$ en una ecuación en diferencias ---una
\emph{relación de recurrencia}--- para la sucesión de coeficientes
$(a_n)_{n\ge0}$. Cambiamos así un problema continuo por uno discreto
(\cref{fig:analogia}). Dos diferencias con respecto a la
\cref{def:diferencias} merecen destacarse desde ahora:
\begin{enumerate}[label=(\roman*)]
  \item los coeficientes de la recurrencia \emph{dependen de $n$}; por
        ejemplo, $y''+y=0$ conduce a $(n+2)(n+1)\,a_{n+2}+a_n=0$;
  \item si los coeficientes de la EDO no son polinómicos, la recurrencia
        involucra \emph{todos} los coeficientes anteriores (es de orden no
        acotado), como veremos en la \cref{sec:recurrencia}.
\end{enumerate}

% Figura: esquema del método (continuo -> discreto -> continuo)
\begin{figure}[ht]
\centering
\begin{tikzpicture}
  \node[nodocian, align=center, text width=4.1cm, minimum height=1.5cm] (edo) at (0,0)
    {\textcolor{cian}{\bfseries EDO}\\[3pt] $P\,y''+A\,y'+B\,y=0$\\[2pt]
     \footnotesize\color{gris} incógnita: la función $y(x)$};
  \node[nodoambar, align=center, text width=4.6cm, minimum height=1.5cm] (rec) at (8.6,0)
    {\textcolor{ambar}{\bfseries Recurrencia}\\[3pt] $a_{n+2}=\Phi_n(a_0,\dots,a_{n+1})$\\[2pt]
     \footnotesize\color{gris} incógnita: la sucesión $(a_n)_{n\ge0}$};
  \node[nodovioleta, align=center, minimum height=1.1cm, inner xsep=10pt] (sol) at (4.3,-2.9)
    {\textcolor{violeta}{\bfseries Solución}\quad
     $y(x)=\sum_{n\ge0}a_n(x-x_0)^n$};

  \draw[flecha, cian] (edo.east) -- (rec.west)
     node[midway, above=3pt, etiqgris] {ansatz en serie}
     node[midway, below=3pt, etiqgris] {continuo $\to$ discreto};
  \draw[flecha, ambar] (rec.south) |- (sol.east)
     node[pos=0.25, right=4pt, etiqgris, align=left] {resolver\\por inducción};
  \draw[flecha, violeta] (sol.west) -| (edo.south)
     node[pos=0.75, left=4pt, etiqgris, align=right] {sumar la serie\\(radio $R>0$)};
\end{tikzpicture}
\caption{Estructura del método de series. La ecuación diferencial para
$y(x)$ se convierte en una ecuación en diferencias para los coeficientes
$a_n$; resuelta ésta, la serie reconstruye la solución en el disco de
convergencia.}
\label{fig:analogia}
\end{figure}

